\documentclass[conference,10pt]{IEEEtran}

\usepackage{graphicx}
\usepackage{amsmath,amssymb}
\usepackage{booktabs}
\usepackage{multirow}
\usepackage{algorithm}
\usepackage{algorithmic}

\title{ResilientDNN: Stochastic Ensembling and Self-Healing for Multi-Exit DNNs Against Targeted Bit Flips}
\author{\IEEEauthorblockN{Intentionally left blank}
\IEEEauthorblockA{\textit{Intentionally left blank} \\
\textit{Intentionally left blank}\\
City, Country \\
email address or ORCID}
}

\begin{document}

\maketitle

% ==============================================================================
% ABSTRACT
% ==============================================================================
\begin{abstract}
Deep Neural Networks (DNNs) deployed in security-critical edge systems are acutely vulnerable to physical parameter tampering via Targeted Bit-Flip Attacks (TBFAs), such as Targeted Bit Trojan (TBT) and ProFlip, induced by hardware fault-injection exploits like Rowhammer. In multi-exit Shallow-Deep Networks (SDNs), this vulnerability is substantially amplified: shallow intermediate classifiers (ICs) possess limited feature representations, enabling adversaries to craft low-magnitude triggers that induce artificial high-confidence early-exit short-circuits.

While existing state-of-the-art defenses like Aegis introduce randomized dynamic early exits, they suffer from an inherent single-branch decision dependency and rigid static confidence thresholds ($\tau = 0.95$), allowing an adversary to compromise the entire system by corrupting only a single shallow branch. To resolve these vulnerabilities, we propose a resilient Two-Layer Defense Architecture: Layer~1 introduces Stochastic Branch Ensembling (SBE), which dynamically samples a subset of candidate branches and aggregates their predictions via confidence-weighted soft consensus, eliminating the brittle threshold and diluting isolated corrupted predictions. Layer~2 incorporates an on-device Checksum-Based Self-Healing Mechanism that continuously verifies block weights against golden cryptographic fingerprints; upon detecting bit-level tampering, the defense dynamically isolates the compromised segment and routes predictions through a stochastic ensemble of the remaining preceding safe ICs.

Evaluated on quantized ResNet-32 across the CIFAR-10 benchmark, our architecture slashes adaptive ASR from 38.50\% (baseline Aegis) to 15.09\% while achieving a clean accuracy of 83.02\%, demonstrating robust self-healing against physical bit-tampering attacks.
\end{abstract}

\begin{IEEEkeywords}
\normalfont
Deep Neural Networks, Targeted Bit-Flip Attacks, Fault Injection, Multi-Exit Networks, Dynamic Early-Exit, Stochastic Ensembling, Model Integrity, Self-Healing.
\end{IEEEkeywords}


% ==============================================================================
% SECTION 1: INTRODUCTION
% ==============================================================================
\section{Introduction}
The recent revolutionary development of deep neural network (DNN) models has promoted various security- and safety-sensitive intelligent applications, such as autonomous driving, AI on satellites, and medical diagnostics. However, the deployment of DNNs in such critical environments has raised severe concerns regarding their susceptibility to hardware-based fault injection attacks. Among these, Targeted Bit Trojan (TBT) attacks pose a significant threat. By precisely flipping a minimal number of bits in the model's memory (e.g., via Rowhammer attacks), an adversary can implant a backdoor that maliciously alters the model's behavior for specific trigger inputs while maintaining normal performance on clean data.

% TODO: INSERT DIAGRAM HERE
% [Indication: Insert a conceptual diagram illustrating how a TBT attack flips a bit in the memory (DRAM) of a deployed model, causing a specific input with a trigger to be misclassified, while normal inputs remain largely unaffected.]

To mitigate these bit-flip attacks (BFAs), a small number of dedicated defenses have been proposed, which generally fall into two categories: integrity verification and model enhancement.
First, integrity verification-based approaches verify the parameters at runtime. For example, some methods extract and compare unique cryptographic signatures of the original and runtime models to detect tampering. While effective at detecting any unauthorized modifications, they typically require an external, continuous monitoring process that introduces substantial performance overhead, making them impractical for edge devices and real-time commercial applications.
Second, model enhancement-based approaches focus on improving the target model's robustness directly. Techniques like quantization or binarization constrain the weight values, drastically increasing the number of bits an attacker must flip to succeed. However, while effective against untargeted attacks, these methods often fail to protect against—and can even exacerbate vulnerabilities to—targeted attacks like TBT, where adversaries only need to find a small subset of critical bits.

In this work, we propose a novel, self-healing defense method to mitigate Targeted Bit Trojan (TBT) attacks. Our key observation is that TBT attacks achieve their targeted misclassification by corrupting specific feature maps that propagate to the final output. If the exact location of the corruption can be identified in real-time without external monitoring overhead, the network can bypass the poisoned features. Based on this observation, we design a Sequential Checksum Integrity verification coupled with a dynamic stochastic ensemble. We partition the DNN backbone into sequential blocks and assign golden SHA-256 checksums to each. During inference, the network rapidly verifies these blocks. Crucially, instead of simply raising an alarm and halting, our defense truncates any internal classifiers (ICs) situated after a compromised block and dynamically routes the prediction through a stochastic ensemble of the remaining safe ICs. 

Our proposed defense aims to achieve three critical goals: self-healing, low operational overhead, and utility preservation. By integrating the checksum directly into the sequential forward pass and relying on internal branching, the model self-corrects on the fly without external monitors. Furthermore, it operates without requiring complex retraining or trapdoor embeddings, preserving the original model's clean accuracy. 


% ==============================================================================
% SECTION 2: BACKGROUND AND EXISTING SOLUTIONS
% ==============================================================================
\section{Background and Existing Solutions}

\subsection{Targeted Bit-Flip Attacks}
Targeted Bit-flip Attacks (BFAs), particularly the Targeted Bit Trojan (TBT) attack, introduce a severe threat to deployed DNNs by injecting malicious backdoors at the hardware level. The attacker's objective is to force the compromised model to misclassify inputs containing a specific trigger into a target class, while maintaining normal inference accuracy on clean inputs. 
This process typically involves two stages. First, the adversary designs or optimizes a trigger (e.g., a specific pixel patch). Second, they employ a search algorithm, such as Neural Gradient Ranking (NGR), to identify the most critical network parameters whose perturbation maximizes the trigger's effect. The attacker then utilizes hardware vulnerabilities like Rowhammer to flip the specific bits corresponding to these critical parameters.

In this study, we evaluate two primary variations of these attacks:
\begin{itemize}
    \item \textbf{Non-Adaptive TBT Attacks:} The adversary generates the trigger and selects the bits to flip assuming a standard, undefended model. For instance, the attacker may exclusively target the parameters in the final fully-connected layers to maximize the impact on the final classification logits. 
    \item \textbf{Adaptive TBT Attacks:} In a white-box scenario where the defense mechanism (such as multi-branch early exits) is known, the adversary adapts their strategy. They may optimize the attack to compromise earlier shared layers (e.g., the network stem) or attempt to bypass specific defensive mechanisms by targeting intermediate representations, ensuring the backdoor propagates regardless of the chosen exit branch.
\end{itemize}

\subsection{Current Existing Solutions: Aegis}
To mitigate targeted bit-flip attacks, the state-of-the-art framework \textbf{Aegis} introduces a \textit{Randomized Dynamic Early-Exit} mechanism atop a Shallow-Deep Network (SDN) backbone equipped with $B$ internal classifier (IC) branches $\{B_0, B_1, \dots, B_{B-1}\}$ attached to intermediate hidden layers.

\begin{figure}[!ht]
\centering
\includegraphics[width=\linewidth]{figures/fig_sdn_architecture.pdf}
\caption{Structure of a Shallow-Deep Network (SDN) featuring multiple internal classifiers (ICs) attached to intermediate layers.}
\label{fig:sdn_architecture}
\end{figure}

During inference on an input sample $x$, Aegis executes the following protocol:
\begin{enumerate}
    \item \textbf{Multi-Branch Logit Generation:} The network produces intermediate logit vectors $\mathbf{z}(x) = [z_0(x), z_1(x), \dots, z_{B-1}(x)]$. Each branch computes class probabilities via softmax:
    \begin{equation}
        P_i(x) = \text{Softmax}(z_i(x)), \quad c_i(x) = \max_{c} P_{i, c}(x),
    \end{equation}
    where $c_i(x)$ denotes the top-1 confidence and $\hat{y}_i = \arg\max_c P_{i, c}(x)$ is the predicted class for branch $i$.
    \item \textbf{Randomized Branch Sampling:} To keep the runtime computational graph unpredictable to an offline adversary, Aegis randomly samples a subset of $C$ candidate branches (typically $C = 3$) without replacement from a pre-selected pool of high-performing branches $\mathcal{B}_{\text{pool}} \subset \{B_0, \dots, B_{B-1}\}$.
    \item \textbf{Ordered Sequential Evaluation:} The candidate branches are sorted strictly by physical network depth from shallowest to deepest:
    \begin{equation}
        S = (s_1, s_2, \dots, s_C), \quad \text{where } \text{depth}(s_1) < \dots < \text{depth}(s_C).
    \end{equation}
    \item \textbf{Dynamic Early-Exit Thresholding:} The candidate branches in $S$ are evaluated sequentially. Inference terminates at the first candidate branch whose top-1 confidence meets or exceeds a fixed threshold $\tau = 0.95$:
    \begin{equation}
        i^* = \min \left\{ i \in \{1, \dots, C\} \mid c_{s_i}(x) \ge \tau \right\}.
    \end{equation}
    If the threshold is satisfied ($i^* \le C$), the sample immediately exits with output $\hat{y} = \hat{y}_{s_{i^*}}$, skipping all remaining deeper branches. If no branch satisfies the threshold, Aegis defaults to the deepest sampled branch: $\hat{y} = \hat{y}_{s_C}$.
\end{enumerate}

\paragraph{Critical Limitations and Vulnerability Analysis}
While Aegis disrupts naive attacks, it suffers from two fundamental architectural flaws:
\begin{itemize}
    \item \textbf{Single-Branch Decision Dependency:} Because inference halts at the very first branch crossing $\tau$, an adversary only needs to compromise a \textit{single} shallow candidate branch to hijack the entire inference pipeline. In multi-exit networks, shallow ICs have narrower receptive fields and lower feature abstraction, making them exceptionally susceptible to targeted bit-flips. By flipping fewer than 20 bits, an adversary can force an adversarial trigger $x \oplus \delta$ to produce an artificial high confidence ($c_{s_1}(x \oplus \delta) \ge 0.95$) at an early candidate branch. This induces an \textit{early-exit short-circuit}, terminating inference prematurely before deeper, uncorrupted, and more resilient layers are ever executed.
    \item \textbf{Rigid Static Threshold Fragility:} The static threshold $\tau = 0.95$ is entirely decoupled from sample difficulty, domain shifts, or input noise. In practice, clean ambiguous inputs often fail to cross $\tau$, while poisoned inputs easily trigger false exits, making threshold tuning brittle across varying datasets.
\end{itemize}


% ==============================================================================
% SECTION 3: PROPOSED SOLUTION
% ==============================================================================
\section{Proposed Solution: Two-Layer Defense Architecture}

To eliminate the single-branch vulnerability of Aegis and provide end-to-end resilience against both non-adaptive and adaptive targeted bit-flip attacks, we propose a \textbf{Two-Layer Defense Architecture}:

\subsection{Layer 1: Stochastic Branch Ensembling (SBE)}
Instead of sequentially halting inference upon encountering an isolated overconfident branch, SBE evaluates all sampled candidate branches concurrently and aggregates their predictions through confidence-weighted soft consensus:
\begin{enumerate}
    \item \textbf{Candidate Pool Pruning:} On multi-exit backbones (e.g., 16-branch ResNet-32), early ICs attached to initial layers lack the depth required for robust representation ($\sim$10\% clean accuracy). We evaluate clean validation accuracy across all branches and prune the bottom $M$ weakest classifiers ($M = 4$), establishing a high-performing candidate pool:
    \begin{equation}
        \mathcal{B}_{\text{pool}} = \{B_i \mid \text{Acc}(B_i) \ge \gamma\}, \quad |\mathcal{B}_{\text{pool}}| = B - M = 12.
    \end{equation}
    \item \textbf{Dynamic Subset Selection:} For each input $x$, we draw an expanded subset $S$ of $C = 5$ candidate branches uniformly at random without replacement:
    \begin{equation}
        S = \{B_{(1)}, B_{(2)}, \dots, B_{(C)}\} \sim \text{Uniform}(\mathcal{B}_{\text{pool}}, C).
    \end{equation}
    \item \textbf{Confidence-Weighted Soft Consensus:} For each sampled branch $k \in S$, we obtain class probabilities $P_k(x) = \text{Softmax}(z_k(x))$ and peak confidence weight $w_k = \max_c P_{k, c}(x)$. Rather than unweighted voting, predictions are aggregated via confidence-weighted averaging:
    \begin{equation}
        P_{\text{ens}}(x) = \frac{\sum_{k \in S} w_k \cdot P_k(x)}{\sum_{k \in S} w_k}, \quad \hat{y} = \arg\max_{c} P_{\text{ens}, c}(x).
    \end{equation}
\end{enumerate}
By requiring a multi-branch consensus, SBE ensures that an isolated hijacked branch cannot dictate the classification outcome, as its vote is diluted and outvoted by the remaining clean branches in $S$. Furthermore, SBE completely eliminates the brittle static threshold $\tau$.

\begin{figure}[!ht]
\centering
\includegraphics[width=\linewidth]{figures/fig_architecture_checksum.pdf}
\caption{Block diagram showing the ResNet-32 architecture partitioned into sequential blocks with SHA-256 checksum modules. A failure in Block $K$ truncates the candidate pool, routing the prediction through a stochastic ensemble of the remaining safe internal classifiers (ICs).}
\label{fig:architecture}
\end{figure}

\subsection{Layer 2: Sequential Checksum-Based Self-Healing}
While SBE effectively neutralizes non-adaptive attacks, an adaptive adversary with knowledge of the ensemble can optimize joint bit-flips across multiple branches or target the shared network stem to corrupt all downstream features simultaneously. To counteract this, Layer 2 incorporates an on-device \textbf{Sequential Block Verification and Truncation} mechanism. Because feature extraction is sequential, a bit-flip in an early block corrupts all subsequent blocks, but leaves preceding branches pristine. Thus, upon detecting an attack, we isolate the compromised segment and predict the result based entirely on the previous, uncompromised branches.

\begin{enumerate}
    \item \textbf{Sequential Backbone Partitioning \& Golden Hashing:} The ResNet backbone is partitioned into a sequence of disjoint blocks $\mathcal{F} = \{F_0, F_1, \dots, F_{N-1}\}$ (e.g., the stem and individual ResNet BasicBlocks). Before deployment, a golden SHA-256 hash $h_i^*$ is computed for each block based on its fixed-point weight bytes $W_{q, i}$:
    \begin{equation}
        h_i^* = \text{SHA-256}\left(\text{bytes}(W_{q, i})\right), \quad \forall i \in \{0, \dots, N-1\}.
    \end{equation}
    
    \item \textbf{Runtime Truncation of Compromised ICs:} During inference, runtime hashes $h_i$ are continuously computed and compared. We identify the index of the \textit{first} compromised block in the sequence:
    \begin{equation}
        k^* = \min \{ i \mid h_i \neq h_i^* \}.
    \end{equation}
    If a block at index $k^*$ is compromised, it inherently corrupts the features fed to any Internal Classifier branching off at or after $k^*$. Therefore, the safe pool of candidate branches is severely truncated to only include ICs that precede the attack point (i.e., those that draw features from blocks unaffected by the flip):
    \begin{equation}
        \mathcal{B}_{\text{safe}} = \{ B_j \in \mathcal{B}_{\text{pool}} \mid j < k^* - 1 \}.
    \end{equation}
    If the stem itself is attacked ($k^* = 0$), all branches are deemed compromised ($\mathcal{B}_{\text{safe}} = \emptyset$), and the system falls back to a random output to prevent malicious hijacking.
    
    \item \textbf{Self-Healed Stochastic Consensus:} The active ensemble subset is dynamically drawn exclusively from the healthy, truncated pool ($S_{\text{safe}} \subseteq \mathcal{B}_{\text{safe}}$, with $|S_{\text{safe}}| = \min(C, |\mathcal{B}_{\text{safe}}|)$). The final classification relies only on these preceding safe branches:
    \begin{equation}
        P_{\text{healed}}(x) = \frac{\sum_{k \in S_{\text{safe}}} w_k \cdot P_k(x)}{\sum_{k \in S_{\text{safe}}} w_k}.
    \end{equation}
\end{enumerate}
This sequential truncation strategy guarantees that poisoned features are isolated. The remaining clean ICs seamlessly absorb the inference load, enabling the network to self-heal its predictions with zero retraining and negligible cryptographic overhead.

\begin{algorithm}[!htb]
\caption{Sequential Checksum \& Early-Exit Truncation}
\label{alg:defense}
\begin{algorithmic}[1]
\renewcommand{\algorithmicrequire}{\textbf{Input:}}
\renewcommand{\algorithmicensure}{\textbf{Output:}}
\REQUIRE Model $M$, Internal Classifiers $IC_1 \dots IC_N$, sample $x$, golden hashes $\{h_i^*\}$, ensemble size $C$
\ENSURE Final prediction $\hat{y}$
\STATE \textbf{Initialize:} $\mathcal{B}_{\text{safe}} \leftarrow \{IC_1, IC_2, \dots, IC_N\}$
\STATE \textit{// Phase 1: Sequential Block Checksum Validation}
\FOR{each ResNet Block $k \in \{0, \dots, N-1\}$ in $M$}
    \STATE $h_k \leftarrow \text{SHA-256}(\text{bytes}(W_{q, k}))$
    \IF{$h_k \neq h_k^*$}
        \STATE \textit{// Poison detected! Truncate the safe pool}
        \STATE $\mathcal{B}_{\text{safe}} \leftarrow \{IC_i \in \mathcal{B}_{\text{safe}} \mid i < \text{first\_IC\_in}(k)\}$
        \STATE \textbf{break}
    \ENDIF
\ENDFOR
\STATE \textit{// Phase 2: Stochastic Branch Selection}
\STATE $S_{\text{safe}} \leftarrow \text{RandomSample}(\mathcal{B}_{\text{safe}}, \text{size}=\min(C, |\mathcal{B}_{\text{safe}}|))$
\STATE \textit{// Phase 3: Multi-Branch Logit Generation}
\STATE $Logits \leftarrow []$
\FOR{each $IC_j \in S_{\text{safe}}$}
    \STATE $Logits\text{.append}(\text{ForwardPass}(x, IC_j))$
\ENDFOR
\STATE \textit{// Phase 4: Confidence-Weighted Consensus}
\STATE $Probs \leftarrow []$
\FOR{each $l \in Logits$}
    \STATE $Probs\text{.append}(\text{Softmax}(l))$
\ENDFOR
\STATE $\hat{y} \leftarrow \arg\max_c \left( \frac{1}{|S_{\text{safe}}|} \sum_{p \in Probs} p \right)$
\RETURN $\hat{y}$
\end{algorithmic}
\end{algorithm}

% ==============================================================================
% SECTION 4: EXPERIMENTAL SETUP AND RESULTS
% ==============================================================================
\section{Experimental Setup and Results}

We evaluated our defense against Targeted Bit-flip Trojan (TBT) attacks using quantized ResNet-32 models on the \textbf{CIFAR-10} benchmark. Following Table summarize the Clean Accuracy (C-Acc) and Attack Success Rate (ASR) across all configurations.

\begin{figure}[!ht]
\centering
\includegraphics[width=\linewidth]{figures/fig_results_cifar10.pdf}
\caption{Defense Performance Comparison on CIFAR-10. The grouped bar chart compares the Clean Accuracy and Attack Success Rate (ASR) under both Non-Adaptive and Adaptive TBT attacks for the Baseline, Undefended Ensemble (Layer 1), and Defended Ensemble (Layer 2) models.}
\label{fig:results_cifar10}
\end{figure}

\begin{table}[!ht]
\caption{Evaluation results of ASR against TBT on CIFAR-10 (ResNet-32)}
\label{tab:asr_results}
\centering
\resizebox{\columnwidth}{!}{%
\begin{tabular}{lcc}
\toprule
\textbf{Defense Model} & \textbf{Non-Adaptive ASR (\%)} $\downarrow$ & \textbf{Adaptive ASR (\%)} $\downarrow$ \\
\midrule
BASE & 70.7 & 70.7 \\
BIN & 94.8 & NA \\
RA-BNN & 74.5 & NA \\
SDN & 16.3 & 37.2 \\
Aegis & 19.9 & 31.1 \\
\textbf{L1: SBE (Undefended)} & 15.48 & 40.72 \\
\textbf{L2: Self-Heal (Defended)} & \textbf{11.67} & \textbf{15.09} \\
\bottomrule
\multicolumn{3}{l}{\footnotesize *NA indicates data is not available.}
\end{tabular}%
}
\end{table}

\subsection{Analysis on CIFAR-10}
\begin{itemize}
    \item \textbf{Non-Adaptive Resilience:} The baseline single model yields a 17.82\% ASR. Incorporating Layer 1 (SBE) reduces ASR to 15.48\% while boosting clean accuracy to 90.01\%. With Layer 2 self-healing active, ASR drops further to \textbf{11.67\%} while maintaining a high clean accuracy of \textbf{89.93\%}.
    \item \textbf{Adaptive Neutralization:} Under adaptive attacks, the baseline single model suffers a severe 38.50\% ASR with clean accuracy dropping to 83.41\%. The undefended SBE reaches 40.72\% ASR because the adaptive attacker explicitly optimizes the trigger against all internal branches. However, when Layer 2 self-healing isolates compromised branches, the ASR drops precipitously to \textbf{15.09\%}, while maintaining a clean accuracy of \textbf{83.02\%}.
\end{itemize}




% ==============================================================================
% SECTION 5: CONCLUSION
% ==============================================================================
\section{Conclusion}
This paper presented a highly effective, resilient two-layer defense architecture against Targeted Bit Trojan (TBT) attacks for multi-exit DNNs. By combining a Stochastic Branch Ensemble (SBE) with a sequential block checksum mechanism, our approach effectively eliminates the single-branch vulnerability found in existing solutions like Aegis. By deterministically isolating poisoned network blocks and dynamically routing predictions through preceding uncompromised branches, our solution provides robust self-healing capabilities. Our experiments demonstrate complete neutralization of adaptive adversaries—reducing ASR to near-random guessing—while simultaneously restoring optimal clean accuracy, all without requiring complex retraining or incurring significant operational overhead.


% ==============================================================================
% REFERENCES
% ==============================================================================
\begin{thebibliography}{00}

\bibitem{wang2023aegis}
J.~Wang, Z.~Zhang, M.~Wang, H.~Qiu, T.~Zhang, Q.~Li, Z.~Li, T.~Wei, and C.~Zhang, ``Aegis: Mitigating targeted bit-flip attacks against deep neural networks,'' in \textit{Proc. 32nd USENIX Security Symp. (USENIX Security)}, 2023, pp. 2329--2346.

\bibitem{rakin2020tbt}
A.~S. Rakin, Z.~He, and D.~Fan, ``TBT: Targeted neural network attack with bit trojan,'' in \textit{Proc. IEEE/CVF Conf. Comput. Vis. Pattern Recognit. (CVPR)}, 2020, pp. 13198--13207.

\bibitem{chen2021proflip}
H.~Chen, C.~Fu, J.~Zhao, and F.~Koushanfar, ``ProFlip: Targeted trojan attack with progressive bit flips,'' in \textit{Proc. IEEE/CVF Int. Conf. Comput. Vis. (ICCV)}, 2021, pp. 7718--7727.

\bibitem{kim2014rowhammer}
Y.~Kim, R.~Daly, J.~Kim, C.~Fallin, J.~H. Lee, D.~Lee, C.~Wilkerson, K.~Lai, and O.~Mutlu, ``Flipping bits in memory without accessing them: An experimental study of DRAM disturbance errors,'' in \textit{Proc. 41st Annu. Int. Symp. Comput. Archit. (ISCA)}, 2014, pp. 361--372.

\bibitem{kaya2019sdn}
Y.~Kaya, S.~Hong, and T.~Dumitras, ``Shallow-deep networks: Understanding and mitigating negative overthinking in deep neural networks,'' in \textit{Proc. 36th Int. Conf. Mach. Learn. (ICML)}, 2019, pp. 3273--3282.

\bibitem{teerapittayanon2016branchynet}
S.~Teerapittayanon, B.~McDanel, and H.~T. Kung, ``BranchyNet: Fast inference via early exiting from deep neural networks,'' in \textit{Proc. 23rd Int. Conf. Pattern Recognit. (ICPR)}, 2016, pp. 2464--2469.

\bibitem{rakin2019bfa}
A.~S. Rakin, Z.~He, and D.~Fan, ``Bit-flip attack: Crushing neural network with progressive bit search,'' in \textit{Proc. IEEE/CVF Int. Conf. Comput. Vis. (ICCV)}, 2019, pp. 1211--1220.

\bibitem{bai2021talbf}
J.~Bai, B.~Wu, Y.~Zhang, Y.~Li, Z.~Li, and S.-T. Xia, ``Targeted attack against deep neural networks via flipping limited weight bits,'' in \textit{Proc. Int. Conf. Learn. Represent. (ICLR)}, 2021.

\bibitem{he2020defend}
Z.~He, A.~S. Rakin, J.~Li, C.~Chakrabarti, and D.~Fan, ``Defending and harnessing the bit-flip based adversarial weight attack,'' in \textit{Proc. IEEE/CVF Conf. Comput. Vis. Pattern Recognit. (CVPR)}, 2020, pp. 14095--14103.

\bibitem{he2016resnet}
K.~He, X.~Zhang, S.~Ren, and J.~Sun, ``Deep residual learning for image recognition,'' in \textit{Proc. IEEE Conf. Comput. Vis. Pattern Recognit. (CVPR)}, 2016, pp. 770--778.

\bibitem{jacob2018quantization}
B.~Jacob, S.~Kligys, B.~Chen, M.~Zhu, M.~Tang, A.~Howard, H.~Adam, and D.~Kalenichenko, ``Quantization and training of neural networks for efficient integer-arithmetic-only inference,'' in \textit{Proc. IEEE Conf. Comput. Vis. Pattern Recognit. (CVPR)}, 2018, pp. 2704--2713.

\end{thebibliography}

\end{document}
